% TEMPLATE-MILC-v3.tex - plantilla maestra UNICA de las guias MILC v3.
%
% La usan las cuatro familias de guias, cada una con su driver:
%   · Grados 6-11          -> scripts/build-guias-g11.py
%   · Bebras               -> scripts/build-guias-bebras.py
%   · Semillero            -> scripts/build-guias-semillero.py
%   · Territorio interior  -> scripts/build-guias-territorio-interior.py
%
% Antes habia tres copias casi identicas (~400 lineas iguales, 6 distintas):
% cada mejora habia que aplicarla tres veces, y en la practica no se hacia
% (el motor de recursos vivia solo en dos de ellas). Lo que varia entre
% familias esta parametrizado en 6 placeholders que cada driver llena
% (se nombran aqui SIN su sintaxis real: el builder reemplaza por texto plano,
% tambien dentro de los comentarios, y un valor multilinea rompe el preambulo):
%
%   HEADER_IZQ        encabezado izquierdo de pagina
%   HEADER_DER        encabezado derecho de pagina
%   PORTADA_ETIQUETA  rotulo sobre el numero grande (GRADO/MOMENTO/MODULO)
%   PORTADA_SERIE     linea de serie de la portada
%   PORTADA_ID        identificador de la guia en la portada
%   PORTADA_CIRCULO   el circulito de grado (°), o vacio si no aplica
%   PDF_TITULO        titulo en texto plano (sin \textbf ni comillas TeX) para
%                     los metadatos del PDF (/Title); lo produce lib_xelatex.texto_pdf
%
% Composicion (auditoria 2026-09-03): las bandas de estacion piden espacio
% (\Needspace*) y prohiben el salto justo despues (\nopagebreak), asi no quedan
% huerfanas al pie; las cajas largas (softbox, drawbox, infoband, puentebox) son
% `breakable`, asi una caja de media pagina ya no arrastra todo a la siguiente
% dejando la anterior al 40 %. Todo texto sobre mostaza o verde va en milcNegro
% (blanco sobre #E5B400 da 1,93:1 y sobre #58A923 2,95:1; ninguno pasa WCAG AA).
% El resto de claves las documenta content/guias/_SCHEMA.md. El script valida
% que no queden placeholders sin reemplazar antes de escribir el archivo final.

% Etiquetado PDF (latex-lab): /Lang, estructura y texto alternativo de las
% figuras, para lectores de pantalla. Debe ir ANTES de \documentclass.
\DocumentMetadata{lang=es-CO,tagging=on}
\documentclass[11pt,letterpaper]{article}
% headheight=24pt: la cabecera derecha lleva dos lineas (identidad de la guia
% y estacion actual). headsep se acorta en la misma medida para que el bloque
% de texto quede exactamente donde estaba con headheight=12pt.
\usepackage[margin=2.0cm,headheight=24pt,headsep=13pt]{geometry}
\usepackage[table]{xcolor}
\usepackage{fontspec}
\usepackage[spanish]{babel}
\usepackage{tikz}
% Librerías TikZ para los diagramas de las guías (bloque `recursos:`):
%  · babel  → babel en español vuelve activos caracteres como «>», y TikZ los usa
%             en sus opciones (>=stealth, ->). Sin esta librería, cualquier
%             diagrama con flechas revienta con «\language@active@arg> has an
%             extra }». Debe ir ANTES de que se dibuje nada.
%  · positioning        → sintaxis «below left=2pt and 3pt».
%  · decorations.pathreplacing → llaves ({brace}) para señalar rangos.
\usetikzlibrary{babel,positioning,decorations.pathreplacing}
\usepackage{graphicx}
% Circuitos: el trabajo de electrónica se hace hoy con micro:bit y sus sensores
% integrados (MakeCode, sin cableado), así que no se necesitan esquemáticos.
% Si algún día entran montajes con protoboard, basta descomentar esta línea:
% los diagramas .tex/.tikz declarados en `recursos:` ya se incrustan solos.
% \usepackage{circuitikz}
\usepackage[most]{tcolorbox}
\usepackage{tabularx}
\usepackage{array}
\usepackage{fancyhdr}
\usepackage{titlesec}
\usepackage[document]{ragged2e}  % bandera a la derecha en todo el cuerpo
\usepackage{enumitem}
\usepackage{microtype}
\usepackage{etoolbox}
\usepackage{needspace}
% hyperref al final del preambulo: metadatos del PDF (titulo, autor, idioma,
% familia), marcadores por estacion (\pdfbookmark) y enlaces sin recuadro.
\usepackage[hidelinks,
  pdftitle={Tecnologías propias — saberes que el mundo aprende ahora},
  pdfauthor={PhD. Álvaro Cárdenas Orozco},
  pdfsubject={Tecnología e Informática · Grado 9 · Guía 9 · MILC v3},
  pdflang={es-CO},
  bookmarksopen=true,bookmarksnumbered=false]{hyperref}

% Helvetica Neue no trae la flecha U+2192, asi que cada «→» escrito literalmente
% en el YAML se compilaba a NADA, en silencio (462 flechas en 61 guias: rutas del
% tipo «paleta → tipografia → lockup» salian rotas). Mapeamos el caracter a su
% version matematica, que si tiene glifo. Asi el autor puede seguir escribiendo
% «→» comodamente en el YAML.
\usepackage{newunicodechar}
\newunicodechar{→}{\ensuremath{\rightarrow}}
% Mismo problema con los simbolos de marca y vinetas: el «✓» de «una palomita ✓»
% salia como cuadrito vacio en el PDF de todas las guias que lo usaban.
\usepackage{pifont}
\newunicodechar{✓}{\ding{51}}
\newunicodechar{✗}{\ding{55}}
\newunicodechar{✔}{\ding{52}}
\newunicodechar{•}{\textbullet}
\newunicodechar{≥}{\ensuremath{\geq}}
\newunicodechar{≤}{\ensuremath{\leq}}

\setmainfont{Helvetica Neue}
\setsansfont{Helvetica Neue}
\setlength{\parindent}{0pt}
\setlength{\parskip}{6pt}
% Sin particion de palabras: en 6.º-7.º una palabra cortada por guion al final
% de linea («Comuni-cacion») frena la lectura mucho mas que una linea corta.
\hyphenpenalty=10000
\exhyphenpenalty=10000
\pretolerance=2000
\tolerance=3000
\setlength{\emergencystretch}{3em}
\linespread{1.10}
\renewcommand{\arraystretch}{1.25}
\setlist{leftmargin=5mm,itemsep=1mm,topsep=1mm}
\newcolumntype{Y}{>{\RaggedRight\arraybackslash}X}

\definecolor{milcMagenta}{HTML}{E6007E}
\definecolor{milcVino}{HTML}{5A0038}
\definecolor{milcTurquesa}{HTML}{0093A5}
\definecolor{milcVerde}{HTML}{58A923}
\definecolor{milcMostaza}{HTML}{E5B400}
\definecolor{milcOcre}{HTML}{B86600}
\definecolor{milcNegro}{HTML}{241816}
\definecolor{milcGris}{HTML}{F7F7F4}
\definecolor{milcLinea}{HTML}{E8E2DD}
\definecolor{milcGradePrimary}{HTML}{7C3AED}
\definecolor{milcGradeDark}{HTML}{4C1D95}
\definecolor{milcGradeSoft}{HTML}{EDE9FE}
\definecolor{milcGradeAccent}{HTML}{CFE7FF}
\definecolor{milcGradeText}{HTML}{FFFFFF}
\color{milcNegro}

\pagestyle{fancy}
\fancyhf{}
% Cabecera de dos lineas: identidad de la guia arriba y, a la derecha y en
% gris, la estacion en curso (\markright desde \milcestacion). Antes de la
% primera estacion la segunda linea va vacia; el \strut mantiene la altura
% para que la linea de arriba no baile entre paginas.
\lhead{\small Tecnología e Informática · Grado 9\\[-1pt]{\footnotesize\strut}}
\rhead{\small Guía 9 · MILC v3\\[-1pt]{\footnotesize\strut\color{milcNegro!65}\rightmark}}
\cfoot{\small\thepage}
\renewcommand{\headrulewidth}{0pt}

% Sobre mostaza (#E5B400) y verde (#58A923) el texto va en milcNegro; sobre el
% resto de la paleta, en blanco. Se usa dentro de fonttitle/fontupper, donde un
% \color posterior manda sobre coltitle.
\newcommand{\milctintasobre}[1]{%
  \ifboolexpr{test{\ifstrequal{#1}{milcMostaza}} or test{\ifstrequal{#1}{milcVerde}}}%
    {\color{milcNegro}}{\color{white}}}

\newtcolorbox{titlebox}[1]{
  enhanced,colback=#1,colframe=#1,arc=7mm,boxrule=0pt,
  left=7mm,right=7mm,top=3mm,bottom=3mm,
  before skip=8pt,after skip=8pt,
  fontupper=\sffamily\bfseries\Large\color{white}
}

\newtcolorbox{softbox}[3]{
  enhanced,breakable,pad at break=3mm,
  colback=#2,colframe=#1,borderline west={4pt}{0pt}{#1},
  arc=2mm,boxrule=.4pt,left=4mm,right=4mm,top=3mm,bottom=3mm,
  before skip=6pt,after skip=8pt,title={#3},coltitle=white,
  colbacktitle=#1,fonttitle=\sffamily\bfseries\milctintasobre{#1},fontupper=\RaggedRight,
  attach boxed title to top left={xshift=3mm,yshift=-2mm},
  boxed title style={arc=2mm, boxrule=0pt}
}

\newtcolorbox{drawbox}[2]{
  enhanced,breakable,pad at break=4mm,
  colback=white,colframe=#1,arc=4mm,boxrule=1pt,
  left=5mm,right=5mm,top=4mm,bottom=4mm,before skip=8pt,after skip=8pt,
  title={#2},coltitle=white,colbacktitle=#1,fonttitle=\sffamily\bfseries\milctintasobre{#1},
  fontupper=\RaggedRight,
  attach boxed title to top left={xshift=4mm,yshift=-2mm},
  boxed title style={arc=3mm, boxrule=0pt}
}

\newtcolorbox{infoband}[1]{
  enhanced,breakable,pad at break=2mm,colback=#1!8,colframe=#1!50,arc=2mm,boxrule=.5pt,
  left=4mm,right=4mm,top=2mm,bottom=2mm,before skip=4pt,after skip=4pt,
  fontupper=\small\RaggedRight
}

% Puente narrativo entre secciones (contrato editorial Conéctate).
% Da continuidad: "Acabas de X. Ahora vas a Y porque Z."
% Visualmente: banda izquierda en color del grado, texto en itálica pequeña.
\newtcolorbox{puentebox}{
  enhanced,breakable,pad at break=2mm,colback=milcGradeSoft,colframe=milcGradeDark,
  borderline west={3pt}{0pt}{milcGradeDark},
  arc=0mm,boxrule=0pt,left=5mm,right=4mm,top=2.5mm,bottom=2.5mm,
  before skip=4pt,after skip=8pt,
  fontupper=\itshape\small\color{milcNegro}\RaggedRight
}
% La flecha va en modo matematico a proposito: Helvetica Neue NO tiene el glifo
% U+2192, asi que \textrightarrow se compilaba a NADA («Missing character: There
% is no → in font Helvetica Neue») y el puente salia sin flecha. En modo
% matematico la flecha viene de la fuente matematica y si se dibuja.
\newcommand{\puente}[1]{%
  \begin{puentebox}%
    \textcolor{milcGradeDark}{$\rightarrow$}\hspace{2mm}#1%
  \end{puentebox}%
}

\newcommand{\linead}[1]{\noindent\color{milcLinea}\rule{#1}{0.5pt}\color{milcNegro}}
\newcommand{\respuesta}[1]{\par\vspace{2mm}\foreach \n in {1,...,#1}{\noindent\color{milcLinea}\rule{\linewidth}{0.5pt}\par\vspace{5mm}}\color{milcNegro}}
\newcommand{\checkbox}{\textcolor{milcGradeDark}{\fbox{\phantom{x}}}\hspace{5pt}}

% ─── Listas aireadas para las guías socioemocionales ────────────────────
% Componer las verificaciones como «\checkbox texto\\» deja el texto que salta
% de línea debajo del cuadrito y apelmaza el bloque. Estos entornos alinean la
% continuación y airean los ítems. Los usa Territorio Interior; cualquier
% familia puede hacerlo.
\newlist{checklista}{itemize}{1}
\setlist[checklista]{
  label=\textcolor{milcGradeDark}{\fbox{\phantom{x}}},
  leftmargin=9mm, labelsep=3.2mm, itemsep=2.4mm,
  topsep=2.5mm, parsep=0pt, partopsep=0pt, before=\RaggedRight
}
\newlist{pasoslista}{enumerate}{1}
\setlist[pasoslista]{
  label=\textcolor{milcGradeDark}{\sffamily\bfseries\arabic*.},
  leftmargin=8mm, labelsep=3mm, itemsep=2.8mm,
  topsep=1mm, parsep=0pt, partopsep=0pt, before=\RaggedRight
}

\newcommand{\phasecard}[4]{
\begin{tcolorbox}[enhanced,colback=#3,colframe=#2,arc=3mm,boxrule=.5pt,height=3.05cm,valign=center,halign=center,left=1.5mm,right=1.5mm,top=2mm,bottom=2mm]
{\sffamily\bfseries\fontsize{22}{24}\selectfont\textcolor{#2}{#1}}\par
{\sffamily\bfseries\small #4}
\end{tcolorbox}
}

% ── Piezas del rediseno 2026 (legibilidad 6.º-11.º) ───────────────────
% Banda de estacion: reemplaza el titlebox macizo. Titulo a la izquierda,
% dato practico (tiempo/modalidad) a la derecha, en una sola linea.
% Nunca huerfana: pide 9 lineas libres (o salta de pagina rellenando con
% \vfill) y prohibe el salto justo despues de la banda. Nueve y no seis:
% la caja partible que sigue necesita ~80pt (titulo adosado, relleno y dos
% lineas) para arrancar en la misma pagina; con menos, tcolorbox la manda
% entera a la siguiente y la banda se queda sola. El penalty 10000 va
% en `after`, ANTES del espacio posterior: un `after skip` (aunque sea 0pt)
% es un \vskip tras la caja, y ese glue es un punto de corte legal que un
% \nopagebreak escrito despues de \end{tcolorbox} ya no cierra. Cada banda
% deja marcador PDF y actualiza la cabecera derecha.
\newcounter{milcestacion}
\newcommand{\milcestacion}[2]{%
\Needspace*{9\baselineskip}%
\stepcounter{milcestacion}%
\pdfbookmark[1]{#2}{milcestacion\themilcestacion}%
\markright{#2}%
\begin{tcolorbox}[enhanced,colback=#1,colframe=#1,arc=2mm,boxrule=0pt,
  left=5mm,right=5mm,top=2.6mm,bottom=2.6mm,before skip=11pt,
  after={\par\nopagebreak\vskip7pt}]
{\sffamily\bfseries\fontsize{15}{18}\selectfont\milctintasobre{#1}#2}
\end{tcolorbox}}

% Tarjeta de paso numerada: sustituye las tablas «Pilar / Aplicacion», que
% eran formato de consulta docente, no de estudiante. El numero va en un
% circulo sobre el borde para que la secuencia se lea de un vistazo.
\newcommand{\milcpaso}[3]{%
\Needspace{4\baselineskip}%
\begin{tcolorbox}[enhanced,colback=white,colframe=milcLinea,arc=1.5mm,boxrule=.9pt,
  left=14mm,right=4mm,top=3mm,bottom=3mm,before skip=4pt,after skip=4pt,
  fontupper=\RaggedRight,
  overlay={\node[anchor=north west,circle,fill=#2,text=white,inner sep=0pt,
    minimum size=7.4mm,font=\sffamily\bfseries\small]
    at ([xshift=3.6mm,yshift=-3.2mm]frame.north west) {#1};}]
#3
\end{tcolorbox}}

% Casilla marcable de tamano util con lapiz (la anterior era \fbox{\phantom{x}},
% de alto variable segun el texto que la seguia).
\newcommand{\milcmarca}{\framebox[3.8mm]{\rule{0pt}{3.4mm}}}

% Fila marcable de lista suelta (checks de la estacion 1).
\newcommand{\milccheckrow}[1]{%
\par\vspace{1.8mm}\noindent
\makebox[6.4mm][l]{\raisebox{-.55mm}{\textcolor{milcNegro}{\milcmarca}}}%
\parbox[t]{\dimexpr\linewidth-6.4mm\relax}{\RaggedRight #1}\par}

% Celda marcable dentro de la rubrica (tabla de autoevaluacion). Misma
% sangria francesa, pero pensada para vivir dentro de una columna X.
\newcommand{\milcopcion}[1]{%
\makebox[5.2mm][l]{\raisebox{-.55mm}{\textcolor{milcNegro}{\milcmarca}}}%
\parbox[t]{\dimexpr\linewidth-5.2mm\relax}{\RaggedRight #1}}

% ── Recursos visuales de la guía ──────────────────────────────────────
% Declarados en el bloque `recursos:` del YAML y emitidos por el builder.
% \guiaFigura   → imagen rasterizada o PDF (foto, carta celeste, montaje).
% \guiaDiagrama → fuente TikZ/circuitikz (.tex) incrustada como vector:
%                 es la vía para circuitos y esquemáticos editables.
% [#1] = texto alternativo (alt) · #2 = ruta relativa al .tex · #3 = pie
% El alt llega al PDF etiquetado (/Alt) y lo lee el lector de pantalla.
\newcommand{\guiaFigura}[3][]{%
\begin{center}
\includegraphics[width=0.88\linewidth,height=0.38\textheight,keepaspectratio,alt={#1}]{#2}\\[1.5mm]
{\footnotesize\itshape\textcolor{milcNegro!75}{#3}}
\end{center}
}

\newcommand{\guiaDiagrama}[2]{%
\begin{center}
\input{#1}\\[1.5mm]
{\footnotesize\itshape\textcolor{milcNegro!75}{#2}}
\end{center}
}

\begin{document}
\thispagestyle{empty}

% ── Portada ───────────────────────────────────────────────────────────
% Antes: un rectangulo de color a pagina completa. Se veia bien en pantalla,
% pero en la fotocopiadora del colegio se comia el toner y salia gris sucio.
% Ahora la banda ocupa el tercio superior y el resto va en blanco.
\begin{tikzpicture}[remember picture,overlay]
  \path[fill=milcGradePrimary]
    (current page.north west) rectangle ([yshift=-10.6cm]current page.north east);

  \node[anchor=north west,text=milcGradeText] at ([xshift=2.0cm,yshift=-1.55cm]current page.north west)
    {\sffamily\bfseries\fontsize{12}{14}\selectfont GRADO};
  \node[anchor=north west,text=milcGradeText,opacity=.85] at ([xshift=2.0cm,yshift=-2.35cm]current page.north west)
    {\sffamily\bfseries\fontsize{8.5}{10}\selectfont SERIE GUÍAS · TECNOLOGÍA E INFORMÁTICA};
  \node[anchor=north east,text=milcGradeText,opacity=.85,align=right] at ([xshift=-2.0cm,yshift=-1.55cm]current page.north east)
    {\sffamily\bfseries\fontsize{9}{12}\selectfont I.E. Sor María Juliana\\Cartago, Valle del Cauca};

  \node[anchor=west,text=milcGradeText] (gradonum) at ([xshift=1.85cm,yshift=-7.1cm]current page.north west)
    {\sffamily\bfseries\fontsize{112}{112}\selectfont 9};
  % El circulito de grado (°) solo aplica a las guias de grado («11°»); en
  % Bebras y Semillero el numero grande es un momento/modulo, no un grado.
  % Se posiciona relativo al nodo `gradonum`, no en coordenadas absolutas.
  \draw[milcGradeText,line width=1.6pt]
    ([xshift=2.0mm,yshift=-4.5mm]gradonum.north east) circle (.15cm);

  \node[anchor=west,text=milcGradeText,text width=11.4cm,align=left] at ([xshift=1.5cm,yshift=0.5cm]gradonum.east)
    {
     {\sffamily\bfseries\fontsize{9}{11}\selectfont\MakeUppercase{Período 1 · Historia de la técnica}}\\[3mm]
     {\sffamily\bfseries\fontsize{21}{25}\selectfont\setlength{\baselineskip}{27pt}Tecnologías propias ---\\[4pt]saberes que el mundo\\[4pt]aprende ahora}\\[3.5mm]
     {\sffamily\bfseries\fontsize{15}{18}\selectfont Guía 9-9-TIC}
    };
\end{tikzpicture}

\vspace*{9.4cm}
\vfill

{\sffamily\bfseries\fontsize{8}{10}\selectfont\textcolor{milcGradeDark}{LO QUE TE LLEVAS HOY}}

\begin{tcolorbox}[enhanced,colback=white,colframe=milcNegro,arc=2mm,boxrule=1.6pt,
  left=5mm,right=5mm,top=4mm,bottom=4mm,before skip=3pt,after skip=12pt,
  fontupper=\RaggedRight\fontsize{11}{15.5}\selectfont]
Una ficha técnica de una página sobre una tecnología propia poco reconocida: andenes andinos, arquitectura en guadua, navegación del Pacífico, trapiche, conuco o la que conozcas de tu región. Lleva el problema técnico que resuelve, el mecanismo dibujado, la evidencia de que funciona, una desventaja honesta frente a la alternativa moderna y un uso actual concreto. Sin un solo argumento que se apoye únicamente en que es tradición.
\end{tcolorbox}

\begin{tcolorbox}[enhanced,colback=milcGris,colframe=milcGris,arc=2mm,boxrule=0pt,
  left=5mm,right=5mm,top=3.5mm,bottom=3.5mm,after skip=12pt]
{\sffamily\bfseries\fontsize{8}{10}\selectfont\textcolor{milcNegro!55}{ESTA GUÍA ES DE}}\\[3mm]
\begin{tabularx}{\linewidth}{X p{3.2cm} p{3.2cm}}
\linead{\linewidth} & \linead{3.0cm} & \linead{3.0cm}\\[-1mm]
{\fontsize{8.5}{10}\selectfont\textcolor{milcNegro!55}{Nombre completo}} &
{\fontsize{8.5}{10}\selectfont\textcolor{milcNegro!55}{Grupo}} &
{\fontsize{8.5}{10}\selectfont\textcolor{milcNegro!55}{Fecha}}\\
\end{tabularx}
\end{tcolorbox}

\vfill

\noindent\textcolor{milcLinea}{\rule{\linewidth}{1pt}}\\[2mm]
{\sffamily\bfseries\fontsize{9.5}{12}\selectfont Autor: Dr. Álvaro Cárdenas Orozco}\\[1mm]
{\sffamily\fontsize{8.5}{11}\selectfont\textcolor{milcNegro!55}{Metodología MILC · apertura ancestral · pensamiento computacional · triángulo Dussel–estoicismo–Floridi}}

\newpage

\section*{Apertura: Tecnologías propias --- saberes que el mundo aprende ahora}

\begin{softbox}{milcGradeDark}{milcGradeSoft}{Saber ancestral}
En los Andes, los agricultores miran las Pléyades en junio para saber si el verano vendrá seco. De eso depende si conviene atrasar la siembra de la papa. Conviene decirlo con precisión: esto es de los Andes de Perú y Bolivia, no del Valle ni de los quimbayas. En el año 2000, tres investigadores lo pusieron a prueba con imágenes de satélite y datos de El Niño, y encontraron que tenían razón. Cuando las Pléyades se ven borrosas es porque hay cirros altos. Y esos cirros anuncian un año de El Niño con menos lluvia (Orlove, Chiang y Cane, 2000). El saber campesino no era superstición: era una observación buena a la que le faltaba la explicación. Hoy en Colombia las Mesas Técnicas Agroclimáticas sientan en la misma mesa a agricultores, meteorólogos y técnicos. Deciden juntos qué sembrar y cuándo (Loboguerrero et al., 2018). La \textbf{cara de exclusión}: el saber local también se equivoca. Presentarlo como infalible es otra forma de faltarle al respeto. El punto no es que los abuelos supieran. Es que una observación se puede poner a prueba.\par\smallskip{\footnotesize\textcolor{milcNegro!70}{\textbf{Fuente.} Orlove, B. S., Chiang, J. C. H., \& Cane, M. A. (2000). Forecasting Andean rainfall and crop yield from the influence of El Niño on Pleiades visibility. \emph{Nature, 403}(6765), 68--71. https://doi.org/10.1038/47456}}
\end{softbox}

\begin{softbox}{milcTurquesa}{milcGris}{Saber contemporáneo}
Una tecnología propia se defiende con cuatro criterios técnicos, no con nostalgia. \textbf{Funcionalidad probada}: resuelve un problema concreto y hay evidencia de que lo resuelve. \textbf{Sostenibilidad}: no agota el recurso del que depende, y a veces lo regenera. \textbf{Adaptación al lugar}: está hecha para ese suelo, ese clima y esos materiales. \textbf{Transmisibilidad}: se puede enseñar sin depender de una fábrica lejana. La prueba de valor no es la antigüedad. Un sistema agrícola que lleva siglos alimentando gente ha pasado más ensayos que cualquier prototipo. Esa es la razón por la que merece atención. Pero una defensa honesta incluye la desventaja. Casi todas estas tecnologías piden más trabajo humano o más conocimiento local que su alternativa industrial. Decirlo no debilita el argumento: es lo que lo vuelve creíble.
\end{softbox}

\begin{softbox}{milcMostaza}{milcGris}{Pregunta-puente}
\textit{Los agricultores andinos leían una señal real y les faltaba la explicación. ¿Qué saber de tu región conoces que nadie ha puesto a prueba todavía, y cómo se pondría a prueba?}
\end{softbox}

\begin{softbox}{milcVerde}{milcGris}{Saber y saber hacer}
Al terminar podrás: (1) \textbf{identificar} el problema técnico que resuelve cada una de cuatro tecnologías propias; (2) \textbf{explicar} los cuatro criterios con que se defiende una tecnología propia sin apelar a la tradición; (3) \textbf{crear} una ficha técnica de una página que defienda una de ellas con evidencia y con una desventaja reconocida.
\end{softbox}



\small
\begin{tabularx}{\textwidth}{>{\bfseries}p{3.0cm}Y}
\rowcolor{milcGradePrimary!12}Autor & Dr. Álvaro Cárdenas Orozco\\
DBA & Análisis crítico de la historia de la tecnología y su impacto social (MEN, T\&I)\\
Referentes & Dussel (técnica situada) · Estoicismo (téchne del cuerpo) · Floridi (infoesfera)\\
Producto final & Una ficha técnica de una página sobre una tecnología propia poco reconocida: andenes andinos, arquitectura en guadua, navegación del Pacífico, trapiche, conuco o la que conozcas de tu región. Lleva el problema técnico que resuelve, el mecanismo dibujado, la evidencia de que funciona, una desventaja honesta frente a la alternativa moderna y un uso actual concreto. Sin un solo argumento que se apoye únicamente en que es tradición.\\
\end{tabularx}
\normalsize

\puente{Vienes de recorrer la historia de la técnica desde el primer filo hasta la era digital. Hoy vuelves la mirada a lo que quedó fuera de ese relato. En la sesión 10 escribes el manifiesto del periodo, y esto es la mitad de su materia.}

\milcestacion{milcGradeDark}{Ruta MILC: así vamos a aprender}
\begin{tabularx}{\textwidth}{>{\centering\arraybackslash}X>{\centering\arraybackslash}X>{\centering\arraybackslash}X>{\centering\arraybackslash}X}
\phasecard{1}{milcVerde}{milcGris}{Escucha\\[-1mm]\small Observo, escucho y pregunto.} &
\phasecard{2}{milcTurquesa}{milcGris}{Entender\\[-1mm]\small Sistematizo y ordeno.} &
\phasecard{3}{milcMagenta}{milcGris}{Hacer\\[-1mm]\small Construyo y aplico.} &
\phasecard{4}{milcVino}{milcGris}{Cierre\\[-1mm]\small Evalúo y reflexiono.}
\end{tabularx}


\puente{Antes de la teoría vas a mirar cuatro tecnologías propias y a preguntarles lo mismo que le preguntarías a una máquina: qué problema resuelve y cómo se sabe que lo resuelve.}

\milcestacion{milcVerde}{Estación 1 de 3 · Escucha}

\begin{softbox}{milcVerde}{milcGris}{Escena para escuchar}
\textbf{Actividad 1 · IDENTIFICA --- Cuatro tecnologías propias} (15 min · individual). (1) Con los materiales del aula o internet, mira estas cuatro: andenes andinos, arquitectura en guadua, navegación del Pacífico y trapiche de caña. (2) Escribe para cada una qué problema técnico resuelve, en una frase. (3) Escribe qué evidencia hay de que lo resuelve: siglos de uso, extensión geográfica o adopción actual. (4) Marca la que más te sorprendió y di en una línea por qué.
\end{softbox}

{\sffamily\bfseries\fontsize{8}{10}\selectfont\textcolor{milcNegro!55}{MARCA CUANDO LO HAYAS HECHO}}
\milccheckrow{Tengo las cuatro con su problema técnico en una frase.}
\milccheckrow{Anoté qué evidencia hay de que cada una funciona.}
\milccheckrow{Marqué la que más me sorprendió y escribí por qué.}

\begin{infoband}{milcVerde}


\textbf{Lo que va al cuaderno (Actividad 1 · IDENTIFICA).} \textbf{Título:} «Actividad 1 --- Cuatro tecnologías propias». \textbf{Formato:} tabla de 4 filas y 3 columnas (tecnología / problema técnico que resuelve / evidencia de que funciona). \textbf{Extensión:} un tercio de página. \textbf{Sabes que terminaste cuando} ninguna casilla de evidencia dice solo «es muy antigua».
\end{infoband}

\puente{Ya tienes las cuatro. Ahora vas a ordenar los criterios con los que se defiende una, para que tu defensa sea técnica y no sentimental.}

\milcestacion{milcTurquesa}{Estación 2 de 3 · Entender}

\begin{softbox}{milcTurquesa}{milcGris}{Lo que vas a entender}
\textbf{Actividad 2 · EXPLICA --- Los cuatro criterios} (20 min · parejas). (1) Con tu pareja, escriban los cuatro criterios con una frase propia cada uno. (2) Apliquen los cuatro a una de las tecnologías de la Actividad 1 y anoten cuáles cumple. (3) Busquen su alternativa industrial y escriban en qué gana cada una. (4) Escriban la desventaja de la tecnología propia sin suavizarla.
\end{softbox}

\milcpaso{1}{milcTurquesa}{\textbf{Los cuatro criterios.} Funcionalidad probada: resuelve un problema y hay evidencia. Sostenibilidad: no agota el recurso del que vive. Adaptación al lugar: está hecha para ese suelo y ese clima. Transmisibilidad: se enseña sin depender de una fábrica lejana.}
\milcpaso{2}{milcTurquesa}{\textbf{La antigüedad no es un argumento, es una pista.} Que algo lleve siglos funcionando indica que pasó muchos ensayos. Pero hay que decir cuáles: qué problema resolvió, dónde, durante cuánto tiempo y con qué resultado medible.}
\milcpaso{3}{milcTurquesa}{\textbf{Toda defensa honesta incluye la desventaja.} Casi todas estas tecnologías piden más trabajo humano o más conocimiento local. Reconocerlo es lo que separa una ficha técnica de un folleto.}
\milcpaso{4}{milcTurquesa}{\textbf{El saber local también se equivoca.} Presentarlo como infalible es otra forma de faltarle al respeto. Lo valioso de las Pléyades no es que los abuelos supieran: es que la observación se pudo poner a prueba y aguantó.}

\begin{drawbox}{milcTurquesa}{Cómo se defiende una tecnología propia}
\textbf{Funcionalidad:} qué problema resuelve y con qué evidencia. \\[1mm] \textbf{Sostenibilidad:} qué recurso usa y si lo agota. \\[1mm] \textbf{Adaptación:} para qué suelo, clima y materiales se hizo. \\[1mm] \textbf{Transmisibilidad:} quién puede aprenderla y con qué. \\[1mm] \textbf{Honestidad:} qué desventaja tiene frente a la alternativa industrial.
\end{drawbox}

\begin{softbox}{milcMostaza}{milcGris}{Errores comunes y su corrección}
(1) \textbf{Defender por antigüedad}: que sea vieja no dice si funciona. (2) \textbf{Ocultar la desventaja}: una ficha sin desventaja se lee como propaganda. (3) \textbf{Presentar el saber local como infalible}: eso también es faltarle al respeto. (4) \textbf{Atribuir un saber al lugar equivocado}: las Pléyades son de los Andes de Perú y Bolivia, no del Valle. (5) \textbf{Describir sin mecanismo}: sin dibujo del cómo funciona, la ficha es una postal.
\end{softbox}

\begin{infoband}{milcTurquesa}


\textbf{Lo que va al cuaderno (Actividad 2 · EXPLICA).} \textbf{Título:} «Actividad 2 --- Los cuatro criterios». \textbf{Formato:} los cuatro criterios con frase propia, aplicados a una tecnología, más la comparación con su alternativa industrial y la desventaja escrita. \textbf{Extensión:} media página. \textbf{Sabes que terminaste cuando} la desventaja está escrita sin suavizar.
\end{infoband}

\puente{Tienes los criterios. Toca elegir una y escribir su ficha, con la desventaja incluida.}

\milcestacion{milcMagenta}{Estación 3 de 3 · Hacer}

\begin{softbox}{milcMagenta}{milcGris}{Cómo vas a construir tu producto}
\textbf{Actividad 3 · CREA --- La ficha técnica} (30 min · individual). (1) Elige una tecnología propia, de las cuatro o de tu región. (2) Escribe el problema técnico que resuelve, en una frase. (3) Dibuja el mecanismo y rotula sus partes: sin dibujo no hay ficha. (4) Escribe la evidencia de que funciona, con lugar y tiempo. (5) Escribe una desventaja honesta frente a la alternativa industrial y un uso actual concreto. (6) Intercambia la ficha con un compañero y marca si algún argumento suyo se sostiene solo en que es tradición.
\end{softbox}

\milcpaso{1}{milcMagenta}{\textbf{Tu ficha tiene cinco partes.} El problema técnico. El mecanismo dibujado y rotulado. La evidencia con lugar y tiempo. La desventaja frente a la alternativa industrial. El uso actual concreto.}
\milcpaso{2}{milcMagenta}{\textbf{El dibujo es obligatorio.} Describir con palabras qué hace un andén es fácil; dibujar cómo retiene el agua y frena la erosión obliga a entenderlo. Ahí se ve quién leyó y quién copió.}
\milcpaso{3}{milcMagenta}{\textbf{La evidencia lleva lugar y tiempo.} «Funciona hace mucho» no sirve. «Se usa en la sierra de Perú desde hace más de mil años, en terrenos con pendiente fuerte» sí sirve.}
\milcpaso{4}{milcMagenta}{\textbf{El uso actual la saca del museo.} Una tecnología con un uso vigente es una propuesta; una sin él es una descripción histórica. Busca dónde se está aplicando hoy y con quién.}

\begin{drawbox}{milcMagenta}{Antes de entregar}
\checkbox El problema técnico está en una frase. \\[1mm] \checkbox Hay dibujo del mecanismo con las partes rotuladas. \\[1mm] \checkbox La evidencia menciona lugar y tiempo. \\[1mm] \checkbox Hay una desventaja escrita sin suavizar. \\[1mm] \checkbox Hay un uso actual concreto, con quién lo aplica. \\[1mm] \checkbox Ningún argumento se sostiene solo en que es tradición.
\end{drawbox}

\begin{softbox}{milcMagenta}{milcGris}{Plantilla de trabajo}
\textbf{Tecnología.} \textit{Cuál y de dónde es, con precisión.} \\[1mm] \textbf{Problema.} \textit{Qué resuelve, en una frase.} \\[1mm] \textbf{Mecanismo.} \textit{Dibujo con partes rotuladas.} \\[1mm] \textbf{Evidencia.} \textit{Lugar, tiempo, resultado.} \\[1mm] \textbf{Desventaja y uso actual.} \textit{Sin suavizar, y con quién la aplica hoy.}
\end{softbox}

\begin{infoband}{milcMagenta}


\textbf{Lo que va al cuaderno (Actividad 3 · CREA).} \textbf{Título:} «Actividad 3 --- La ficha técnica». \textbf{Formato:} las cinco partes de la ficha, con el dibujo del mecanismo rotulado y la desventaja escrita. \textbf{Extensión:} una página. \textbf{Sabes que terminaste cuando} tu ficha se sostiene sin usar ni una vez la palabra «ancestral» como argumento.
\end{infoband}

\puente{Lo que entregas hoy: una ficha de una página con mecanismo dibujado y evidencia. La prueba de calidad: no queda ni un argumento que se sostenga solo en que es tradición.}

\milcestacion{milcMostaza}{Producto de la sesión}
\textbf{Ficha técnica de una página con mecanismo dibujado, evidencia y desventaja reconocida}.
\begin{softbox}{milcMostaza}{milcGris}{Criterios de calidad}
(1) El problema técnico está enunciado en una frase. (2) Hay dibujo del mecanismo con las partes rotuladas. (3) La evidencia menciona lugar y tiempo, no solo antigüedad. (4) Hay una desventaja honesta frente a la alternativa industrial. (5) Hay un uso actual concreto y se dice quién lo aplica. (6) Revisaste la ficha de un compañero y marcaste sus argumentos de sola tradición.
\end{softbox}

\puente{Antes de cerrar, mira lo que hiciste desde cinco lados. Defender un saber sin ponerlo a prueba es tratarlo como adorno, no como conocimiento.}

\milcestacion{milcVino}{Evaluación liberadora · cinco dimensiones}

\begin{softbox}{milcVino}{milcGris}{Sentido de la evaluación}
Evaluar no es solo revisar una nota. Es reconocer qué aprendí, qué cambió en mí, qué emociones manejé, cómo me relaciono con la ciudad y la comunidad, y qué le dejaré a quien viene detrás.
\end{softbox}

\textbf{Mi reflexión liberadora en cinco dimensiones.} Completa cada frase iniciada:

\small
\begin{tabularx}{\textwidth}{>{\bfseries\RaggedRight\arraybackslash}p{3.4cm}Y>{\RaggedRight\arraybackslash}p{4.6cm}}
\rowcolor{milcVino}\color{white}Dimensión & \color{white}\bfseries Mi respuesta (frame starter) & \color{white}\bfseries Algo que descubrí\\
1. Personal & Lo que más cambió en mí fue \linead{6.5cm} & \linead{4.2cm}\\
2. Emocional & La emoción más fuerte fue \linead{2.5cm} y la regulé \linead{3cm} & \linead{4.2cm}\\
3. Ciudadana & En democracia esto importa porque \linead{6cm} & \linead{4.2cm}\\
4. Local (Cartago) & En mi barrio sería útil para \linead{6.5cm} & \linead{4.2cm}\\
5. Intergeneracional & A mi abuela/abuelo le diría \linead{6.5cm} & \linead{4.2cm}\\
\end{tabularx}
\normalsize

\begin{infoband}{milcVino}
\textbf{Para la próxima sustentación.} Vuelve a esta tabla en seis meses. Verás que las respuestas cambian — no porque lo aprendido cambie, sino porque tú cambias. La evaluación liberadora es una conversación contigo a través del tiempo.
\end{infoband}


\puente{Para terminar, tres citas verificadas sobre saber y probar: Dussel sobre la tecnología humanizada, Marco Aurelio sobre mirar el curso de los astros, Floridi sobre los patrones pequeños que solo significan si se comparan.}

\milcestacion{milcGradeDark}{Triángulo de pensamiento · cierre filosófico}

\begin{softbox}{milcVino}{milcGris}{Dussel · lente del nosotros}
\textit{«No hay liberación sin economía y tecnología humanizada, diseño, y sin partir de una formación social histórica.»}

{\footnotesize\itshape\color{milcNegro!70}--- Enrique Dussel · Filosofía de la liberación (1977), §2.6.7.1}

Una tecnología propia parte de su lugar: ese suelo, ese clima, esa gente. No es una versión pobre de otra industrial. Es una respuesta distinta a un problema que la industrial resolvió pensando en otro sitio.

\textbf{¿Qué problema de mi región se está resolviendo con una solución pensada para otro lugar?}
\end{softbox}
\linead{\linewidth}\\[1mm]
\linead{\linewidth}

\begin{softbox}{milcMostaza}{milcGris}{Estoicismo · Marco Aurelio · Meditaciones VII, 47 (c. 175 d.C.) · cuidado interior}
\textit{«Conduce mirar alrededor el curso de los astros, como quien gira con ellos, y contemplar también frecuentemente las mutuas conversiones de los elementos, porque las consideraciones de estas cosas purifican a uno de las manchas de esta vida terrestre.»}

{\footnotesize\itshape\color{milcNegro!70}--- Marco Aurelio · Meditaciones VII, 47 (c. 175 d.C.)}

Mirar el cielo para decidir cuándo sembrar es exactamente eso: leer el curso de los astros y actuar. Los agricultores andinos lo hacían con las Pléyades, y en el año 2000 un satélite confirmó que la señal existía.

\textbf{¿Qué señal del entorno miro sin darme cuenta y uso para decidir algo?}
\end{softbox}
\linead{\linewidth}\\[1mm]
\linead{\linewidth}

\begin{softbox}{milcTurquesa}{milcGris}{Floridi · lente de la infoesfera}
\textit{«Los pequeños patrones solo pueden ser significativos si se agregan correctamente, se comparan y se procesan a tiempo. (trad. propia, abreviada)»}

{\footnotesize\itshape\color{milcNegro!70}--- Luciano Floridi · Big data and their epistemological challenge (2012)}

El brillo de las Pléyades una noche no dice nada. Cientos de observaciones a lo largo de generaciones, comparadas y usadas a tiempo, sí dicen algo. Eso es lo que un satélite pudo verificar después.

\textbf{¿Qué observación mía repetida muchas veces podría significar algo si la comparara?}
\end{softbox}
\linead{\linewidth}\\[1mm]
\linead{\linewidth}

\begin{infoband}{milcNegro}
\textbf{Nota para el docente.} Las tres son citas textuales y llevan su referencia debajo. Puedes remitir a la obra en clase.
\end{infoband}

\milcestacion{milcVino}{Mi autoevaluación · marca una casilla por fila}

{\small
\setlength{\extrarowheight}{2.2pt}
\begin{tabularx}{\textwidth}{>{\RaggedRight\arraybackslash\bfseries}p{3.0cm}YYY}
\rowcolor{milcVino}\color{white}\bfseries Criterio & \color{white}\bfseries Lo logré & \color{white}\bfseries En proceso & \color{white}\bfseries Necesito apoyo\\[.6mm]
Comprensión del tema & \milcopcion{Explico las ideas con mis palabras.} & \milcopcion{Reconozco las ideas, falta precisión.} & \milcopcion{Necesito repasar lo básico.}\\[.8mm]
\rowcolor{milcGris}Pensamiento computacional & \milcopcion{Usé los pilares al armar mi producto.} & \milcopcion{Usé algunos pilares.} & \milcopcion{Necesito releer la estación 2.}\\[.8mm]
Aplicación y producto & \milcopcion{Mi producto cumple los criterios.} & \milcopcion{Está armado, con ajustes pendientes.} & \milcopcion{Aún está en construcción.}\\[.8mm]
\rowcolor{milcGris}Reflexión 5 dimensiones & \milcopcion{Respondí cada una con honestidad.} & \milcopcion{Respondí algunas con detalle.} & \milcopcion{Mis respuestas son muy generales.}\\[.8mm]
Triángulo de pensamiento & \milcopcion{Conecté las 3 voces con mi trabajo.} & \milcopcion{Conecté al menos una.} & \milcopcion{No las relaciono con mi proceso.}\\[.8mm]
\end{tabularx}}

\vspace{3mm}
\textbf{Hoy entendí que:}\\[1mm]
\linead{\linewidth}\\[1mm]
\linead{\linewidth}

\textbf{Mi compromiso para la próxima sesión es:}\\[1mm]
\linead{\linewidth}\\[1mm]
\linead{\linewidth}

\begin{softbox}{milcMostaza}{milcGris}{Cita de despedida · Marco Aurelio}
\textit{``El obstáculo en el camino se vuelve el camino. Lo que estorba al camino se vuelve camino.''}\\[1mm]
Cada error de hoy, bien observado, se convierte en pieza del camino que pisarás mañana.
\end{softbox}

\small
\begin{tabularx}{\textwidth}{>{\bfseries}p{3.6cm}Y}
\rowcolor{milcGradeDark!12}Nombre completo: & \linead{10cm}\\
Fecha: & \linead{4cm} \quad Curso/grupo: \linead{4cm}\\
Mi firma: & \linead{9cm}\\
\end{tabularx}
\normalsize

\begin{infoband}{milcGradeDark}
\textbf{Para la próxima sesión.} Llega con la ficha completa y el dibujo del mecanismo. Pregúntale a alguien mayor de tu familia por una técnica de su oficio que no se enseñe en ningún curso. En la sesión 10 escribes el manifiesto del periodo, y esa respuesta te va a servir.
\end{infoband}

\end{document}
